\documentclass[11pt]{article}
\usepackage{icmlko}
\begin{document}

\title{천장을 열한 중 열에서 넘었다 --- 모형 쪽에 남은 것은 $+0.0009$, 자료 쪽은 그 서른한 배}
\author{월드모델 실험실}
\date{노트 505 · 2026년 8월 3일}
\maketitle

\begin{abstract}
노트 490 이 자료 손잡이 하나에 $-0.0283$ 을 매기고 노트 $497 \sim 504$ 가
여덟 노트를 들여 모형에서 $+0.0042$ 를 벌었다. 그러면 물음은 하나다 ---
\textbf{모형 쪽에 아직 얼마가 남았나.}

먼저 노트 454 의 자(오라클 칸 상한)를 열한 도메인에 대 봤다. \textbf{퇴화
한다} --- 여덟 도메인에서 행/칸 $= 1.00$ 이라 칸이 곧 행이고 상한이
정의상 $+1.0000$ 이다. 소진율은 그냥 챔피언 점수를 다시 쓴 것이 된다.

\begin{quote}
\textbf{조항} --- \textbf{오라클 칸 상한은 행/칸 $> 1$ 일 때만 자다.}
이 자료에서 쓸 수 있는 곳은 시장팝업($8.40$) · 만화($1.39$) ·
팝업($1.16$) 셋뿐이다.
\end{quote}

그래서 \textbf{같은 시기 천장}(그 도메인 유보 행만으로 $5$ 겹
교차검증 --- 노트 $444 \cdot 445$ 의 자)으로 바꿔 쟀다.

\begin{center}
\begin{tabular}{lrrr}
\hline
도메인 & 같은 시기 천장 & 챔피언 & 비 \\
\hline
만화 & $+0.1679$ & $+0.3364$ & $\mathbf{2.00}$ \\
아이돌 & $+0.2326$ & $+0.4442$ & $\mathbf{1.91}$ \\
펀딩 & $+0.2109$ & $+0.3634$ & $\mathbf{1.72}$ \\
팝업 & $+0.2382$ & $+0.4088$ & $\mathbf{1.72}$ \\
모바일 & $+0.5614$ & $+0.6223$ & $1.11$ \\
세계애니 & $+0.5355$ & $+0.5727$ & $1.07$ \\
애니 & $+0.5613$ & $+0.5792$ & $1.03$ \\
게임 & $+0.6519$ & $+0.6624$ & $1.02$ \\
웹툰 & $+0.4730$ & $+0.4769$ & $1.01$ \\
시장팝업 & $-0.1177$ & $+0.1481$ & (천장 음수) \\
\hline
도서 & $+0.3890$ & $+0.3698$ & $0.95$ \\
\hline
\end{tabular}
\end{center}

\textbf{열한 중 열에서 챔피언이 그 도메인의 자기 천장을 넘는다.} 여지가
남은 곳은 도서 하나이고, 판 무게를 곱하면 \textbf{$+0.0009$} 다.
\end{abstract}

\begin{figure}[t]
\centering
\includegraphics[width=\textwidth]{figs/f1.pdf}
\caption{\textbf{가} --- 천장 $\to$ 챔피언. \textbf{나} --- 크기와
넘는 배율. \textbf{다} --- 남은 곳.}
\end{figure}

\section{작을수록 많이 넘는다}

넘는 배율이 판 무게와 뚜렷이 반대다 --- 아이돌($1.5\%$) $1.91$ · 팝업
($1.9\%$) $1.72$ · 만화($7.5\%$) $2.00$ · 펀딩($15.4\%$) $1.72$ 대
웹툰($20.7\%$) $1.01$ · 애니($17.7\%$) $1.03$ · 게임($5.2\%$) $1.02$.

\textbf{노트 445 의 ``모음은 작은 도메인만을 위한 것''이 전면 확인된다.}
거기서는 여섯 도메인의 부분 측정이었고 여기서는 열한 전부다. 큰 도메인은
자기 자료로 이미 거의 다 하고, 모으기는 작은 도메인을 끌어올린다.

\textbf{정직하게 --- 이 천장은 약한 천장이다.} 한 도메인의 유보 행
$51 \sim 711$ 개로 $5$ 겹을 돌린 것이라 상한이라기보다 ``자기 자료만으로는
여기까지''다. \textbf{넘었다는 것은 모형이 최적이라는 뜻이 아니라
모으기가 실제로 일하고 있다는 뜻이다.} 그러나 \emph{모형 개선의 여지}를
재는 데는 맞는 자다 --- 더 짜낼 신호가 그 도메인 안에 있다면 그 도메인
자기 모형이 먼저 찾아냈을 것이다.

\section{그래서 어디에 남았나}

\begin{center}
\begin{tabular}{lr}
\hline
& 판 크기 \\
\hline
모형 쪽 남은 여지(천장까지) & $+0.0009$ \\
자료 손잡이 하나(검색 시계열, 노트 490) & $\mathbf{0.0283}$ \\
\hline
비 & $\mathbf{31}$ 배 \\
\hline
\end{tabular}
\end{center}

그리고 그 자료는 지금 \textbf{덮음 $41.7\%$} 에서만 일한다 --- 노트 491
이 찾은 \texttt{[:20]} 절단 때문에 묶여 있었다.

\begin{quote}
\textbf{이 사이클의 결론} --- 여덟 노트가 모형에서 $+0.0042$ 를 벌었고,
같은 기간에 잰 결과 \textbf{모형 쪽에 남은 것은 $+0.0009$} 다.
\textbf{모형 갈래는 사실상 닫혔다.} 다음 배율은 수집에서만 나온다.
\end{quote}

\section*{남은 것}

\begin{itemize}
\item \textbf{도서 $+0.0192$} --- 유일하게 천장 아래인 도메인. 작지만
유일하다.
\item \textbf{시장팝업 천장이 음수다}($-0.1177$) --- 자기 자료로는 아예
못 배운다. 챔피언 $+0.1481$ 은 \textbf{전부 남에게서 빌린 것}(노트 445 ·
453 의 그 도메인).
\item \textbf{자료 쪽 세 갈래} --- ① 검색 시계열 재수집(절단 고침 뒤,
$+152$ 행 예보) ② 시장팝업 태거(칸 $15$ 개) ③ CN 태거(축 셋).
\textbf{전부 사람이 정할 일이다.}
\end{itemize}

\end{document}
